\chapter{Campo dei reali non standard}\stepcounter{capitoloeccolo}

Si parla di modelli non standard dei numeri reali quando si considerano campi reali \emph{non archimedei}, che verifichino un certo principio di \straniero{transfer}: in questo modo le proprietà elementari dei numeri reali sono conservate. Inoltre, in questi campi, esistono numeri infinitesimi ed infiniti, che si comportano come utilizzati da Lebniz. 

In questo capitolo si costruirà formalmente un campo $ ^{*}\R $ di numeri reali non standard, e se ne vedranno le principali proprietà.

\definizione{
    Un campo ordinato $ \mathcal{K}=(\K, +,-, \cdot ,0,1,\le) $ si dice \emph{archimedeo} se per ogni $ x \in \K $, con $ 0<x $, esiste $ n \in \N $ tale che l'elemento \[\parentesi{n\text{-volte}}{1+\dots+1} \eqqcolon n1\]sia $ x<n1 $.
}

% % \section{Esistenza di un campo non archimedeo}

% \definizione{
%     Un campo ordinato $ \mathcal{K}=(\K, +,-, \cdot ,0,1,\le) $ si dice \emph{archimedeo} se per ogni $ x \in \K $, con $ 0<x $, esiste $ n \in \N $ tale che l'elemento \[\parentesi{n\text{-volte}}{1+\dots+1} \eqqcolon n1\]sia $ x<n1 $
% }
% \teorema{ijpijoijodijoijoijodijcoijoijoijoijoijoij}{
%     Per ogni campo ordinato archimedeo esiste un campo ordinato non archimedeo elementarmente equivalente al primo.    
% % Ogni campo ordinato è elementarmente equivalente ad un campo ordinato non Archimedeo. 
% }
% \dimostrazione{ijpijoijodijoijoijodijcoijoijoijoijoijoij}{
%     Si estenda il linguaggio dei campi: \[
%         \mathcal{L}_{\text{Fld}}=\{+,-, \cdot ,0,1,\le\} \subseteq \mathcal{L}'=\mathcal{L}_{\text{Fld}}\cup\{c\}
%     \]dove $ c $ è una costante, e si consideri una $\mathcal{L}_{\text{Fld}}$-struttura $ \mathcal{K}$. 

%     Sia ora $ T $ una $ \mathcal{L}'$-teoria:\[
%         T=\operatorname{Th}_{{\text{Fld}}}(\mathcal{K})\cup\{\tau_{n}: 1\le n \in \N \}
%     \]dove $ \tau_{n}  $ è l'enunciato $ n \cdot 1 \le c $.

%     $ T $ è finitamente soddisfacibile. Infatti, sia $ 1\le n \in \N $. Allora la $ \mathcal{L}' $-espasione $ \mathcal{K}_{n}  $ tale che $ c^{\mathcal{K}_{n} } = n \cdot 1 $ è tale che \[
%         \mathcal{K}_n \vDash \operatorname{Th}_{{\text{Fld}}}(\mathcal{K})\cup\{\tau_{i}: 1\le i \le n, i \in \N \}.
%     \]
    
%     Per il teorema di compattezza esiste $ \mathcal{F}' $, una $ \mathcal{L}' $-struttura tale che \[
%         \mathcal{F}'\vDash T'.
%     \]Se $ \mathcal{F} $ è la $ \mathcal{L}_{\text{Fld}} $-ridotta di $ \mathcal{F}' $, \[
%         \mathcal{F}=\left(A,+,-, \cdot ,0,1,\le\right)
%     \]allora in $ \mathcal{F} $ l'elemento $ c^{\mathcal{F}'} \in A $ è tale che $ n \cdot 1\le  c^{\mathcal{F}'} $ per ogni $ n\ge 1 $. 
    
%     $\implies$ $ \mathcal{F} $ non è Archimedeo e $ \mathcal{F} $ è elementarmente equivalente a $ \mathcal{K} $.\qed
% }

% \osservazione{
%     Il teorema \teoref{ijpijoijodijoijoijodijcoijoijoijoijoijoij} assicura che esista un campo $ \starR $ elementarmente equivalente ad $ \R $ ma non archimedeo. 
% }

\section{Costruzione degli iperreali}

L'obiettivo di questa sezione è costruire $ \starR $ come ultraprodotto di $ \mathcal{L}$-strutture.

\subsection{Prodotto di $ \mathcal{L} $-strutture}

\paragrafo{Prodotto Cartesiano}{%
    Sia $ (A_{i} )_{i \in I} $ una famiglia di insiemi non vuoti indicizzata su $ I\neq \emptyset $. Il \emph{prodotto cartesiano} degli $ A_{i}$ è l'insieme%
    \[
        \prod_{i \in I} A_{i}=\left\{f:I\longrightarrow \bigcup_{i \in I} A_{i}: f(i) \in A_{i}\text{ per ogni } i \in I  \right\}  
    \]Gli elementi $ f \in \prod_{i \in I} A_{i} $ si indicano come sequenze generalizzate $ \sequence{a_{i}}_{i \in I} $ dove $ a_{i}=f(i)  $ è la $ i $-esima coordinata di $ \sequence{a_{i} }_{i \in I} $ (o di $ f $) per ogni $ i \in I $. 

    Quando è chiaro l'insieme su cui varia $ i $, si userà $ \sequence{a_{i}} $ in luogo di $ \sequence{a_{i}}_{i \in I} $.

    Quando $ A_{i} =A $ per ogni $ i \in I $ si scriverà $ A^{I} $ invece di $ \prod_{i \in I} A_{i} $.
}{}{}
\paragrafo{Prodotto di strutture}{%
    Sia $ (\mathcal{A}_i)_{i \in I} $ una famiglia di $ \mathcal{L} $-strutture (con dominio, rispettivamente, $ A_{i}  $) indicizzata su $ I\neq \emptyset $. Il loro \emph{prodotto (diretto)} $ \prod_{i \in I} \mathcal{A}_{i} $ è la $ \mathcal{L} $-struttura $ \mathcal{M} $ con dominio $ \prod_{i \in I} A_{i} $ e simboli di $ \mathcal{L} $ interpretati come segue: \begin{itemize}
        \item per ogni $ R \in \operatorname{Rel} $ con arietà $ k $, ed ogni $ \sequence{a_{i}^{1}},\dots,\sequence{a_{i}^{k}} \in\prod_{i \in I} A_{i}$ \[
            \left(\sequence{a_{i}^{1}},\dots,\sequence{a_{i}^{k}}\right) \in R^{\mathcal{M}}\qquad \text{se e solo se}\qquad \left(a_{i}^{1}, \dots, a_{i}^{k}  \right) \in R^{\mathcal{A}_i}\text{ per ogni }i \in I;
        \]
        \item per ogni $ f \in \operatorname{Fun} $ con arietà $ k $, ed ogni $ \sequence{a_{i}^{1}},\dots,\sequence{a_{i}^{k}}, \sequence{b_{i} } \in\prod_{i \in I} A_{i}$ \[
            f^{\mathcal{M}}\left(\sequence{a_{i}^{1}},\dots,\sequence{a_{i}^{k}}\right) = \sequence{b_{i} }\qquad \text{se e solo se}\qquad \parbox{12em}{$f^{\mathcal{A}_i}\left(a_{i}^{1}, \dots, a_{i}^{k}  \right) = b_{i} $ per ogni $i \in I$;}
        \]
        \item per ogni $ c \in \operatorname{Cost} $ ed ogni $ \sequence{b_{i}} \in \prod_{i \in I}A_{i}   $ \[
            c^{\mathcal{M}}=\sequence{b_{i}}\qquad\text{se e solo se}\qquad c^{\mathcal{A}_i}=b_{i} \text{ per ogni }i \in I.
        \]
    \end{itemize}
}{}{}
\subsection{Filtri e ultrafiltri}

\definizione{
    Sia $ I\neq \emptyset $ un insieme. Una famiglia non vuota $ \mathcal{F} \subseteq \parti{I} $ è un \emph{filtro} se soddisfa le seguenti proprietà;\begin{itemize}
        \item se $ A \in \mathcal{F} $ e $ A \subseteq B \subseteq I $, allora $ B \in \mathcal{F} $;
        \item se $ A,B \in \mathcal{F} $, allora $ A \cap B \in \mathcal{F}$.
    \end{itemize}

    Un filtro $ \mathcal{F} $ si dice \emph{proprio} se $ \emptyset\notin \mathcal{F} $.
}
\definizione{
    Il \emph{filtro di Fréchet} $ \mathcal{F}_{\text{Fr}}^{I} $ su $ I $ è definito nel seguente modo:\citazione{$ A \in \mathcal{F}_{\text{Fr}}^{I} $ se e solo se $ I\setminus A $ è finito.}
}
\paragrafo{FIP}{%
    Una famiglia $ \mathcal{G} \subseteq\parti{I} \setminus\{\emptyset\}$ ha la \emph{proprietà dell'intersezione finita} (brevemente FIP) se ogni intersezione finita di elementi di $ \mathcal{G} $ è non vuota.
}{}{}
\paragrafo{Estensione canonica}{%
    Ogni famiglia $ \mathcal{G}\subseteq\parti{I} \setminus\{\emptyset\}$ con la FIP si estende canonicamente al filtro proprio\[
        \langle \mathcal{G}\rangle =\{ B \subseteq I: B \supseteq A_1\cap\dots\cap A_{k}\text{ per qualche }k \in \N\text{ e }A_1,\dots,A_{k} \in \mathcal{G}\}.
    \]Equivalentemente, $ \langle \mathcal{G}\rangle $ è il più piccolo filtro su $ I $ che contiene $ G $
}{}{}
\paragrafo{Osservazione}{%
    $ \langle \mathcal{G}\rangle $ è un filtro proprio. Infatti:\begin{itemize}
        \item se $ A \in \langle \mathcal{G}\rangle $, allora per qualche $ A_1,\dots,A_{k} \in \mathcal{G}  $ \[
            A \supseteq A_1 \cap \dots\cap A_{k} 
        \]e se $ B \supseteq A $ allora \[
            B \supseteq A_1 \cap \dots\cap A_{k} 
        \]e quindi $ B \in \langle \mathcal{G}\rangle $; 
        \item se $ B_1, B_2 \in \langle \mathcal{G}\rangle $, allora \begin{align*}
            B_1 &\supseteq A_1\cap \dots\cap A_{k}\\
            B_2 &\supseteq C_1\cap \dots\cap C_{k}
        \end{align*}e allora \[
            B_1\cap B_2 \supseteq A_1\cap \dots\cap A_{k}\cap C_1\cap \dots\cap C_{k}
        \]e quindi $ B_1\cap B_2 \in \langle \mathcal{G}\rangle$. 
    \end{itemize}

    La condizione FIP di $ \mathcal{G} $ garantisce che $ \langle \mathcal{G}\rangle $ sia un filtro \emph{proprio}, ma è possibile costruire un filtro in questo modo a partire da una qualsiasi famiglia di sottoinsiemi di $ I $. 
}{}{}
\definizione{
    Un filtro proprio $ \mathcal{F} $ su $ I $ si dice \begin{itemize}
        \item \emph{ultrafiltro} se per ogni $ A \subseteq I $ si ha $ A \in \mathcal{F} $ o $ I\setminus A \in \mathcal{F} $; 
        \item \emph{primo} se $ A\cup B \in \mathcal{F} $ implica che $ A \in \mathcal{F} $ o $ B \in \mathcal{F} $;
        \item \emph{massimale} se non esiste un filtro proprio $ \mathcal{F}' $ su $ I $ tale che $ \mathcal{F}' \supsetneqq \mathcal{F} $.
    \end{itemize}
}
\proposizione{doijcoijoosijoijoijdcidicoijcoijdoijoijoij}{
    Dato $ I \neq \emptyset$, $ \mathcal{F} $ filtro proprio su $ I $, allora sono fatti equivalenti: \begin{enumerate}
        \item $ \mathcal{F} $ è un ultrafiltro;
        \item $ \mathcal{F} $ è primo;
        \item $ \mathcal{F} $ è massimale.
    \end{enumerate} 
}
\dimostrazioneprop{doijcoijoosijoijoijdcidicoijcoijdoijoijoij}{
    $ \boxed{1.\Rightarrow 2.} $\hspace{1em} Per assurdo, $ A\cup B \in \mathcal{F} $, ma $ A, B\notin \mathcal{F} $. 
    
    Per $ 1. $ si ha che $ I \setminus A, I \setminus B \in \mathcal{F} $. Allora\[
        \parentesi{=I\setminus (A\cup B)}{(I\setminus A) \cap (I \setminus B)} \in \mathcal{F}
    \]che è un assurdo perché $ \mathcal{F} $ è proprio.

    $ \boxed{2.\Rightarrow 3.} $\hspace{1em} Si supponga per assurdo che esista un filtro proprio $ \mathcal{F}'\supsetneqq \mathcal{F} $. Sia $ A \in \mathcal{F}' \setminus \mathcal{F} $. 
    
    In particolare, $ I \setminus A \notin \mathcal{F}' $, quindi $ I \setminus A \notin \mathcal{F}$. 
    
    Ma $ A \cup (I\setminus A) = I \in \mathcal{F} $, e $ A, I\setminus A \notin \mathcal{F} $, il che contraddice 2.

    $ \boxed{3.\Rightarrow 1.} $\hspace{1em} Si supponga per assurdo che $ A \notin \mathcal{F} $ e $ I\setminus A \notin \mathcal{F} $. 
    
    Sia $ \mathcal{G}=\mathcal{F}\cup \{A\} $. $ \mathcal{G} $ ha la FIP\footnote{Infatti, siano $ A_1,\dots,A_{k} \in \mathcal{G}   $, senza perdita di generalità assunti tutti diversi. Se $ A_1,\dots,A_{k} \in \mathcal{F}  $, allora $ A_1\cap\dots\cap A_{k}\neq \emptyset  $ perché $ \mathcal{F} $ è un filtro proprio. Se invece $ A_{k} = A  $, allora \[
        \parentesi{ \in \mathcal{F}}{A_1\cap \dots\cap A_{k-1} }\cap A =  B \cap A
    \]con $B = A_1\cap \dots\cap A_{k-1}$. Se capitasse che $ B\cap A =\emptyset$, allora $ B \subseteq I\setminus A $, da cui $ I\setminus A \in \mathcal{F} $, assurdo poiché per ipotesi $ I \setminus A \notin \mathcal{F} $.}. Sia ora $ \mathcal{F}'=\langle \mathcal{G}\rangle $. $ \mathcal{F}' $ è un filtro proprio, e si ha \[
        \mathcal{F}' \supseteq \mathcal{G}\supseteq \mathcal{F}, \qquad \mathcal{F}'\neq \mathcal{F}
    \]poiché $ A \in \mathcal{F}'\setminus \mathcal{F} $. Questo contraddice 3.\qed
}
\definizione{
    Un ultrafiltro $ \mathcal{U} $ su $ I $ si dice \emph{non principale} se per ogni $ A \subseteq I $ finito, $ I\setminus A \in \mathcal{U} $.
}
\paragrafo{Lemma degli ultrafiltri}{%
    Ogni filtro proprio $ \mathcal{F} $ su $ I $ può essere esteso ad un ultrafiltro $ \mathcal{U}\supseteq \mathcal{F} $ su $ I $.
}{doijcsoijoijoijoijoijoijoij}{}
\dimostrazioneprop{doijcsoijoijoijoijoijoijoij}{
    Sia $ X=\{\mathcal{G}: \mathcal{G} \supseteq \mathcal{F}, \mathcal{G}\text{ filtro proprio su }I\} $. 

    Si consideri $ (X, \subseteq)$: sia $ C \subseteq X $ una catena, cioè se $ \mathcal{G}, \mathcal{G}' \subseteq C $ allora o $ \mathcal{G} \subseteq \mathcal{G}' $ oppure $ \mathcal{G}' \subseteq \mathcal{G} $. 

    Sia $ \displaystyle \mathcal{H} = \bigcup_{\mathcal{G} \in C} \mathcal{G} $. 
    
    È evidente che per ogni $ \mathcal{G} \in C $, $ \mathcal{G} \subseteq \mathcal{H} $. Necessariamente $\emptyset \notin \mathcal{H}$ perché $ \mathcal{H} $ unione di filtri propri, e inoltre $ \mathcal{H} \supseteq \mathcal{F} $ poiché unione di filtri che contengono $ \mathcal{F} $. Inoltre $ \mathcal{H} $ è un filtro. Infatti: \begin{itemize}
        \item \emph{chiusura per sovrainsiemi}: se $ A \in \mathcal{H} $, allora $ A \in \mathcal{G} $ per qualche $ \mathcal{G} \in C $ dunque, per ogni $ B \supseteq A $ necessariamente $ B \in \mathcal{G} $, poiché $ \mathcal{G} $ è filtro, e dunque $ B \in \mathcal{H} $;
        \item \emph{intersezioni finite}: siano $ A, B \in \mathcal{H} $: si ha che $ A \in \mathcal{G} $ e $ B \in \mathcal{G}' $, per qualche $ \mathcal{G}, \mathcal{G}'  \in C$; poiché $ C $ è una catena, allora necessariamente deve valere una delle seguenti: \[
            A, B \in \mathcal{G} \supseteq \mathcal{G}',\qquad A, B \in \mathcal{G}' \supseteq \mathcal{G};
        \]e quindi $ A\cap B \in \mathcal{G} $ (oppure $ A\cap B \in \mathcal{G}' $) poiché $ \mathcal{G} $ (oppure $ \mathcal{G}' $) è un filtro. Dunque $ A\cap B \in \mathcal{H} $.
    \end{itemize}

    Per il lemma di Zorn, esiste $ \mathcal{U} $ massimale in $ X $, e necessariamente $ \mathcal{U} \supseteq \mathcal{F} $. Dunque $ \mathcal{U} $ è un filtro massimale rispetto alla famiglia di tutti i filtri su $ I $\footnote{Se per assurdo $ \mathcal{U} $ non fosse un ultrafiltro massimale, sia $ \mathcal{U}' \supsetneqq \mathcal{U} $ un filtro proprio su $ I $. Dunque $ \mathcal{U}' \supseteq \mathcal{F}$ e $ \mathcal{U}' $ filtro proprio su $ I $, e quindi $ \mathcal{U}' \in X $. Questo contraddice il risultato secondo cui $ \mathcal{U} $ è massimale in $ X $.}, e per la proposizione \framref{doijcoijoosijoijoijdcidicoijcoijdoijoijoij}, un filtro massimale è un ultrafiltro.\qed
}
\subsection{Ultraprodotto di $ \mathcal{L} $-strutture}

\definizione{
    Sia $ (A_{i} )_{i \in I} $ una famiglia di insiemi non vuota indicizzata da $ I \neq \emptyset $. Dato un filtro $ \mathcal{F} $ su $ I $ si definisce la relazione di equivalenza $ \sim_{\mathcal{F}}  $ su $ \prod_{i \in I} A_{i}   $ ponendo per $ \sequence{a_i},\sequence{b_i} \in \prod_{i \in I} A_{i} $ \[
        \sequence{a_i}\sim_{\mathcal{F}}\sequence{b_i}\qquad\text{se e solo se}\qquad \{i \in I:a_i=b_i \} \in \mathcal{F}. 
    \]La classe di equivalenza rispetto a $ \sim_{\mathcal{F}}$ di un elemento $ \sequence{a_i} \in \prod_{i \in I}A_{i}$ verrà indicata con $ [\sequence{a_i}] $
}
\osservazione{
    $ \sim_{\mathcal{F}} 
    $ è una relazione di equivalenza: simmetria e riflessività sono ovvie; per la transitività: siano $ \sequence{a_i}\sim_{\mathcal{F}} \sequence{b_i} $, e $ \sequence{b_i}\sim_{\mathcal{F}}\sequence{c_i} $. Allora \begin{align*}
            X_{\sequence{a_i},\sequence{b_i}}&\coloneqq\{i \in I: a_i = b_i\} \in \mathcal{F}\\
            X_{\sequence{b_i},\sequence{c_i}}&\coloneqq\{i \in I: b_i = c_i\} \in \mathcal{F}
        \end{align*}

    Si deve dimostrare che l'insieme $ X_{\sequence{a_i},\sequence{c_i}} \coloneqq\{i \in I: a_i = b_i\} \in \mathcal{F}$, ma \[
        X_{\sequence{a_i},\sequence{b_i}} \cap X_{\sequence{b_i},\sequence{c_i}} \subseteq X_{\sequence{a_i},\sequence{c_i}}   
    \]ma $ X_{\sequence{a_i},\sequence{b_i}} \cap X_{\sequence{b_i},\sequence{c_i}} \in \mathcal{F} $, e dunque $ X_{\sequence{a_i},\sequence{c_i}} \in \mathcal{F} $.
}
\paragrafo{Prodotto ridotto di insiemi}{%
    Il \emph{prodotto ridotto} $ \prod_{{\mathcal{F}}} A_{i}$ degli $ A_{i}  $ modulo $ \mathcal{F} $ è il quoziente \[
        \prod_{i \in I}A_{i}/\sim_{\mathcal{F}} = \left\{\left[\sequence{a_{i}}\right]: \sequence{a_{i}} \in \prod_{i \in I}A_{i}  \right\}. 
    \]Se $ \mathcal{F} $ è un ultrafiltro, il prodotto ridotto corrispondente $ \prod_{\mathcal{F}}A_{i}   $ si dice \emph{ultraprodotto}. Quando $ A_{i}=A  $ per ogni $ i \in I $, si parla invece di \emph{potenza ridotta} e scriviamo $ A^{I}/\mathcal{F} $; se $ \mathcal{F} $ è un ultrafiltro, si parla di \emph{ultrapotenza}.
}{}{}
\paragrafo{Ultraprodotto di strutture}{%
    Sia $ (\mathcal{A}_i)_{i \in I} $ una famiglia di $ \mathcal{L} $-strutture indicizzata da $ I\neq \emptyset $, e sia $ \mathcal{U} $ un ultrafiltro su $ I $. L'ultraprodotto $ \prod_{\mathcal{U}}\mathcal{A}_i  $ è la $ \mathcal{L} $-struttura $ \mathcal{M} $ con dominio $ \prod_{\mathcal{U}}A_{i}   $ e simboli di $ \mathcal{L} $ interpretati come segue: \begin{itemize}
        \item per ogni $ R \in \operatorname{Rel} $ con arietà $ k $ ed ogni $ \sequence{a_{i}^{1}},\dots,\sequence{a_{i}^{k}} \in \prod_{i \in I}A_{i}  $%
            \begin{gather*}
                \left(\left[\sequence{a_{i}^{1}}\right],\dots,\left[\sequence{a_{i}^{k}}\right]\right) \in R^{\mathcal{M}}\\ 
                \text{se e solo se}\\ 
                \left\{i \in I: \left(a_{i}^{1}, \dots, a_{i}^{k}  \right) \in R^{\mathcal{A}_i}\right\} \in \mathcal{U};
            \end{gather*}
        \item per ogni $ f \in \operatorname{Fun} $ con arietà $ k $ ed ogni $ \sequence{a_{i}^{1}},\dots,\sequence{a_{i}^{k}} ,\sequence{b_{i} } \in \prod_{i \in I}A_{i}  $%
            \begin{gather*}
                f^{\mathcal{M}}\left(\left[\sequence{a_{i}^{1}}\right],\dots,\left[\sequence{a_{i}^{k}}\right]\right)=\left[\sequence{b_{i} }\right]\\ 
                \text{se e solo se}\\ 
                \left\{i \in I: f^{\mathcal{A}_i}\left(a_{i}^{1}, \dots, a_{i}^{k}  \right) = b_{i} \right\} \in \mathcal{U};
            \end{gather*}
        \item per ogni $ c \in \operatorname{Cost} $ ed ogni $ \sequence{b_{i} } \in \prod_{i \in I}A_{i}   $ \[
            c^{\mathcal{M}}=\left[\sequence{b_{i} }\right]\qquad\text{se e solo se}\qquad \{i \in I: c^{\mathcal{A}_i}=b_{i} \} \in \mathcal{U}.
        \]
    \end{itemize}
    Se $ \mathcal{A}_i = \mathcal{A}$ per ogni $  i \in I $, la $ \mathcal{L} $-struttura $ \mathcal{A}^{I}/\mathcal{U} $ si dice \emph{ultrapotenza} di $ \mathcal{A} $.
}{dcjkklkmdjjjjultraprodottoinssimo}{}
\osservazione{
    Le definizioni sopra riportate sono ben date. A titolo esplicativo si dimostra che se $ R $ è un simbolo di relazione binario, e $ \sequence{a^{1}_i}, \sequence{a^{2}_i },\sequence{b^{1}_i}, \sequence{b^{2}_i } \in \prod_{i \in I} A_{i}  $ con \begin{equation*}
        \sequence{a_{i}^{1} } \sim_{\mathcal{U}} \sequence{b_{i}^{1} } ,\qquad \sequence{a_{i}^{2} } \sim_{\mathcal{U}} \sequence{b_{i}^{2} } 
    \end{equation*}allora $  \left\{i \in I: \left(a_{i}^{1}, a_{i}^{2}  \right) \in R^{\mathcal{A}_i}\right\} \in \mathcal{U} $ se e solo se $  \left\{i \in I: \left(b_{i}^{1}, b_{i}^{2}  \right) \in R^{\mathcal{A}_i}\right\} \in \mathcal{U} $. Per simmetria si dimostra un'unica implicazione.

    Per definizione, i seguenti insiemi appartengono ad $ \mathcal{U} $: \begin{equation*}
        \set{i \in I: a^{1}_i = b^{1}_i}, \quad\set{i \in I: a^{2}_i = b^{2}_i}. 
    \end{equation*}Inoltre vale \begin{equation*}
        \set{i \in I: \left(a_{i}^{1}, a_{i}^{2}  \right) \in R^{\mathcal{A}_i}}\cap\set{i \in I: a^{1}_i = b^{1}_i}\cap\set{i \in I: a^{2}_i = b^{2}_i} \subseteq \set{i \in I: \left(b_{i}^{1}, b_{i}^{2}  \right) \in R^{\mathcal{A}_i}}
    \end{equation*}infatti, fissato $ i \in I $, se $ \left(a_{i}^{1}, a_{i}^{2}  \right) \in R^{\mathcal{A}_i} $ e $ a_{i}^{1} = b_{i}^{1}   $ e $ a_{i}^{2} = b^{2}_i  $, allora $ \left(b_{i}^{1},b_{i}^{2}  \right) = \left(a_{i}^{1}, a_{i}^{2}  \right) \in R^{\mathcal{A}_i}$.

    Questo dimostra la buona definizione della definizione precedente nel caso in una relazione binaria.
}
\proposizione{doijcoijdijcijijoijoijoijcoaijoijoijoij}{
    Sia $ (\mathcal{A}_i)_{i \in I} $ una famiglia di $ \mathcal{L} $-strutture (con dominio, rispettivamente, $ A_{i}  $) indicizzata su $ I \neq \emptyset $, e sia $ \mathcal{U} $ un ultrafiltro su $ I $. Per ogni $ \mathcal{L} $-termine $ t(x_1,\dots,x_{n} ) $ ed ogni $ \sequence{a_{i}^{1}},\dots,\sequence{a_{i}^{n}} ,\sequence{b_{i} } \in \prod_{i \in I}A_{i}   $\begin{gather*}
        t^{\mathcal{M}}\left(\left[\sequence{a_{i}^{1}}\right],\dots,\left[\sequence{a_{i}^{k}}\right]\right)=\left[\sequence{b_{i} }\right]\\ 
        \text{se e solo se}\\ 
        \left\{i \in I: t^{\mathcal{A}_i}\left(a_{i}^{1}, \dots, a_{i}^{k}  \right) = b_{i} \right\} \in \mathcal{U};
    \end{gather*}
}
\dimframmento{doijcoijdijcijijoijoijoijcoaijoijoijoij}{
    Si dimostra per induzione sulla complessità del termine. \begin{itemize}
        \item \emph{Caso base:} i termini di complessità minima sono le costanti e le variabili. Nel caso in cui il termine $ t $ sia una costante, il risultato è valido per definizione di ultraprodotto. 
        
        Quanto $ t(x): x $ è una costante, fissati $ \sequence{a_{i} },\sequence{b_{i} } \in \prod_{i \in I} A_{i}   $, \begin{equation*}
            t^{\mathcal{M}}\left(\left[\sequence{a_{i} }\right]\right) = \left[\sequence{b_{i} }\right]
        \end{equation*}se e solo se $ \left[\sequence{a_{i} }\right]=\left[\sequence{b_{i} }\right] $, ovvero, per definizione di classi di equivalenze, se e solo se \begin{equation*}
            \set{i \in I: t^{\mathcal{A}_i}(a_{i} ) = b_{i} }=\set{i \in I: a_{i} = b_{i}  } \in \mathcal{U}.
        \end{equation*}
        \item \emph{Passo induttivo:} siano $ t_1(x_1,\dots,x_{n} ), \dots, t_k(x_1,\dots,x_{n} ) $ termini per cui valga il risultato, e sia $ f $ simbolo di funzione di arietà $ k $. 
        
        Sia dunque $ t(x_1,\dots,x_{n}) $ il termine \begin{equation*}
            f\left(t_1(x_1,\dots,x_{n} ), \dots, t_{k}(x_1,\dots,x_{n} ) \right).
        \end{equation*} \commento{Non sono riuscito a terminare questa dimostrazione}
        % Fissati $ \sequence{a^{1}_i},\dots,\sequence{a^{n}_i}, \sequence{b_{i} } \in \prod_{i \in I} A_{i}  $, si dimostrerà che \begin{equation*}
        %     t^{\mathcal{M}}\left(\left[\sequence{a^{1}_i}\right],\dots,\left[\sequence{a^{n}_i}\right]\right) = \left[\sequence{t^{\mathcal{A}_i}\left(a^{1}_i,\dots,a^{n}_i\right)}\right].
        % \end{equation*}Questo corrisponde alla tesi.
        % % \begin{equation*}
        % %     t_{j}^{\mathcal{M}}\left(\left[\sequence{a^{1}_i}\right],\dots,\left[\sequence{a^{n}_i}\right]\right) = \left[\sequence{c^{j}_i}\right]
        % % \end{equation*}per qualche $ \sequence{c^{j}_i} \in \prod_{i \in I} A_{i}   $.

        % Dunque $ t^{\mathcal{M}} \left(\left[\sequence{a^{1}_i}\right],\dots,\left[\sequence{a^{n}_i}\right]\right) = \left[\sequence{b_i}\right]$ se e solo se \begin{equation*}
        %     f^{\mathcal{M}}\left(t_{1}^{\mathcal{M}}\left(\left[\sequence{a^{1}_i}\right],\dots,\left[\sequence{a^{n}_i}\right]\right),\dots,t_{k}^{\mathcal{M}}\left(\left[\sequence{a^{1}_i}\right],\dots,\left[\sequence{a^{n}_i}\right]\right)\right) = \left[\sequence{b_{i} }\right]
        % \end{equation*}e a sua volta questo è equivalente ad \begin{equation}
        %     f^{\mathcal{M}}\left(\left[\sequence{c^{1}_i}\right],\dots,\left[\sequence{c^{k}_i}\right]\right) = \left[\sequence{b_{i} }\right].\label{eq:ciaonen}
        % \end{equation}Per la definizione di funzioni in un ultraprodotto, \eqref{eq:ciaonen} se e solo se \begin{equation*}
        %     B \coloneqq \set{i \in I: f^{\mathcal{A}_i}\left(c^{1}_i,\dots,c^{k}_i\right)=b_{i} } \in \mathcal{U}.
        % \end{equation*}Si considerino ora gli insiemi \begin{align*}
        %     B' &\coloneqq\set{i \in I: t^{\mathcal{A}_i}(a^{1}_i,\dots,a^{n}_i) = b_{i} },\\ 
        %     B''&\coloneqq B\cap \set{i \in I: t_{1}^{\mathcal{A}_i}(a^{1}_i,\dots,a^{n}_i) = c^{1}_i }\cap \dots \cap \set{i \in I: t_{k}^{\mathcal{A}_i}(a^{1}_i,\dots,a^{n}_i) = c^{k}_i }
        % \end{align*}L'inclusione $ B'' \subseteq B' $ è banale, e pertanto se $ B \in \mathcal{U} $ allora $ B' \in \mathcal{U}$, poiché $ \mathcal{U} $ è un filtro.

        % Viceversa, se $ B' \in \mathcal{U}$,
    \end{itemize}
}
\teorema[Teorema di Łoś]{dsoijsocijoijdididcoijdoijoijoicjoijoijocijoijoijcoijoijoij}{
    Sia $ (\mathcal{A}_i)_{i \in I} $ una famiglia di $ \mathcal{L} $-strutture (con dominio, rispettivamente, $ A_{i}  $) indicizzata su $ I \neq \emptyset $, e sia $ \mathcal{U} $ un ultrafiltro su $ I $. Sia $ \mathcal{M}\coloneqq\prod_{\mathcal{U}} \mathcal{A}_{i}   $. Per ogni $ \mathcal{L} $-formula $ \varphi(x_1,\dots,x_{n} ) $ e per ogni $\sequence{a_{i}^{1}},\dots,\sequence{a_{i}^{n}} \in\prod_{i \in I} A_{i}$, si ha \[
        \mathcal{M}\vDash \varphi\left[\left[\sequence{a_{i}^{1} }\right],\dots,\left[\sequence{a_{i}^{n} }\right]\right]\quad\text{se e solo se}\quad X_{\varphi, \sequence{a_{i}^{1}},\dots,\sequence{a_{i}^{n}}} \in \mathcal{U}
    \]dove $ X_{\varphi, \sequence{a_{i}^{1}},\dots,\sequence{a_{i}^{n}}} = \left\{i \in I: A_{i}\vDash \varphi\left[a_{i}^{1},\dots,a_{i}^{n}  \right] \right\} $
}
\dimostrazione{dsoijsocijoijdididcoijdoijoijoicjoijoijocijoijoijcoijoijoij}{
    Si dimostra il teorema per induzione sull'altezza della formula. \begin{itemize}
        \item \emph{Caso base:} per dimostrare il caso di una $ \mathcal{L}$-formula atomica, per semplicità, si considererà un simbolo di relazione binaria $ R $ e la formula $ \varphi: R(t_1,t_2) $. 
        
        Siano dunque $ \left[\sequence{b^{1}_i}\right] = t_1^{\mathcal{M}},\quad \left[\sequence{b^{2}_i}\right] = t_2^{\mathcal{M}} $. Necessariamente, per definizione di termine in un ultraprodotto, \begin{align*}
            X_{t_1}&=\{i \in I: t_1^{\mathcal{A}_i} = b_{i}^{1} \} \in \mathcal{U},\\
            X_{t_2}&=\{i \in I: t_2^{\mathcal{A}_i} = b^{2}_i\} \in \mathcal{U}.
        \end{align*}
        
        Per definizione, $ \mathcal{M}\vDash R(t_1,t_2) $ se e solo se $ \left[\sequence{b^{1}_i}\right] R^{\mathcal{M}} \left[\sequence{b^{2}_i}\right] $ e, per la definizione di relazione in un ultraprodotto, questo si verifica se e solo se $ X_{R} \in \mathcal{U}  $, dove \[
            X_{R}=\{i \in I: b^{1}_i\, R^{\mathcal{\mathcal{A}}_i}\, b^{2}_i\}.
        \]

        % Si ha che $ X_{R} \in \mathcal{U}  $ sse 
        Sia ora\[
            X=\{i \in I: b^{1}_i\, R^{\mathcal{\mathcal{A}}_i}\, b^{2}_i\,\text{ e }\,t_1^{\mathcal{A}_i} = b_{i}^{1}\,\text{ e }\,t_2^{\mathcal{A}_i} = b^{2}_i\}.
        \]$X \in \mathcal{U}$ se e solo se $ X_{R} \in \mathcal{U} $, poiché $ X_{R}\cap X_{t_1} \cap X_{t_2} \subseteq X$ e $ \mathcal{U} $ è un ultrafiltro.

        Ma $ \displaystyle X = \{i \in I: \mathcal{A}_i\vDash R(t_1,t_2)\} = X_{\varphi} $, e si è quindi dimostrata la tesi.
        \item \emph{Ipotesi induttiva:} l'ipotesi induttiva si verifica unicamente per i connettori e quantificatori logici $ \lnot, \lor, \exists$, in quanto tutti i rimanenti sono esprimibili come combinazione di questi.
        \begin{itemize}
            \item \emph{Negazione:} sia $ \varphi $ la formula $ \lnot\,\psi $, tale che il teorema valga per $ \psi $. 
            
            Per definizione $ \mathcal{M}\vDash \varphi $ se e solo se $ \mathcal{M}\nvDash \psi $, e, per ipotesi induttiva, questo è valido se e solo se $ \quad X_{\psi}\notin \mathcal{U}  $.
            
            Ma $ X_{\psi}\notin\mathcal{U}  $ se e solo se $ X_{\varphi} \in \mathcal{U}  $, poiché $ X_{\varphi}=I\setminus X_{\psi} $ e $ \mathcal{U} $ è ultrafiltro.
            \item \emph{Disgiunzione:} sia $ \varphi $ la formula $ \psi \,\lor\, \chi $, tale per cui il teorema valga per $ \psi $ e $ \chi $. 
            
            Per definizione $ \mathcal{M}\vDash \varphi $ se e solo se $ \mathcal{M}\vDash \psi $ oppure $ \mathcal{M}\vDash \chi $, e, per ipotesi induttiva, questo è equivalente a 
            \[
                X_{\psi} \in \mathcal{U}\,\text{ o }\,X_{\chi} \in \mathcal{U}  
            \]Inoltre, poiché $ \mathcal{U} $ è un filtro primo (condizione equivalente ad essere ultrafiltro) \[
                X_{\psi} \in U\,\text{ o }\,X_{\chi} \in \mathcal{U} \quad\text{se e solo se}\quad X_{\psi}\cup X_{\chi} \in \mathcal{U}   
            \]e, visto che $ X_{\varphi} = X_{\psi \,\lor\,\chi} = X_{\psi}\cup X_{\chi}     $, questo dimostra la tesi.
            \item \emph{Esistenziale:} sia $ \varphi $ la formula $ \exists y \, \psi(y) $, tale per cui il teorema valga per $ \psi(y) $. 
            
            $ \mathcal{M}\vDash \exists y \, \psi(y)$ se e solo se esiste $ \sequence{b_{i} } \in \prod_{i \in I} A_{i} $ tale per cui $ \mathcal{M}\vDash \psi\left[\left[\sequence{b_{i} }\right]\right] $. 
            
            Se esiste $\sequence{b_{i} } \in \prod_{i \in I} A_{i} $ tale per cui $ \mathcal{M}\vDash \psi\left[\left[\sequence{b_{i} }\right]\right]$, allora $ X_{\psi, \sequence{b_{i}}} \in \mathcal{U} $ per ipotesi induttiva. Inoltre, $ X_{\psi, \sequence{b_{i} }} \subseteq X_{\varphi} $ poiché \begin{align*}
                X_{\psi,\sequence{b_i}} &= \{ i \in I: \mathcal{A}_i\vDash \psi\left[b_{i} \right]\}\\ 
                X_{\varphi} &= \{i \in I: \mathcal{A}_i\vDash \exists\, y\,\psi\} 
            \end{align*}e quindi $ X_{\varphi} \in \mathcal{U}$ poiché $ \mathcal{U} $ è un ultrafiltro.
            
            Viceversa, se $ X_{\varphi} \in \mathcal{U}$, per ogni $ i \in  X_{\varphi}  $ sia un $ a_{i} \in A_{i}$ fissato tale che $ \mathcal{A}_i \vDash \psi[a_{i} ]$ (utilizzando l'assioma di scelta). Per ogni $ i \notin X_{\varphi}  $, sia $ c_{i} \in A_{i}$ scelto in maniera arbitraria. 

            Si definisce ora $ \sequence{b_{i} } \in \prod_{i \in I}A_{i}$: \[
                b_{i}  = \begin{cases}
                    a_{i} & i \in X_{\varphi}\\ 
                    c_{i} & i \notin X_{\varphi}     
                \end{cases} 
            \]e si ha che, necessariamente, \[
                X_{\psi,\sequence{b_{i}}} = \left\{i \in I: \mathcal{A}_i\vDash \psi\left[b_{i} \right]\right\} \supseteq X_{\varphi} \in \mathcal{U}
            \]e dunque $ X_{\psi,\sequence{b_{i}}} \in \mathcal{U}  $. Per ipotesi induttiva, $ \mathcal{M}\vDash\psi\left[\left[\sequence{b_{i} }\right]\right] $, e quindi $ \mathcal{M}\vDash \exists y\,\psi(y) $.\qed
        \end{itemize}
    \end{itemize}
}
\paragrafo{Corollario}{%
    Sia $ \mathcal{A} $ una $ \mathcal{L} $-struttura e $ \mathcal{U} $ un ultrafiltro su $ I $. Sia $ \pi: \mathcal{A}\longrightarrow \mathcal{A}^{I}/\mathcal{U} $ definita da $ \pi(a)=[\sequence{a_{i} }] $ dove $ a_{i} \equiv a $ per ogni $ i \in I $. Allora $ \pi $ è un'immersione elementare e, in particolare, $ \mathcal{A} $ e $ \mathcal{A}^{I}/\mathcal{U} $ sono elementarmente equivalenti.
}{dkjnckjjjjjjjjjjjjjdjjjjjdjnkjnkjn}{}
\dimframmento{dkjnckjjjjjjjjjjjjjdjjjjjdjnkjnkjn}{
    $ \pi $ è un morfismo. Questo segue banalmente dalla definizione degli elementi di $ \mathcal{A}^{I}/\mathcal{U} $, tenendo in considerazione che $ I \in \mathcal{U} $. 

    Si indichi con $ A $ il dominio di $ \mathcal{A} $, e con $ A^{I}/\mathcal{U} $ il dominio di $ \mathcal{A}^{I}/\mathcal{U} $.

    Per dimostrare che $ \pi $ sia una immersione elementare si sfrutta la proposizione \framref{odijcoijodijocijoijoijoi}. Si vuole dimostrare che per ogni $ \mathcal{L} $-formula $ \phi(y,x_1,\dots,x_{n}) $ e per ogni $ a^{1},\dots,a^{n} \in A $, \begin{gather}
        \mathcal{A}^{I}/\mathcal{U}\vDash \exists y \,\phi\left[y,\pi(a^{1}), \dots, \pi(a^{n})\right]\label{eq:uno}\\ 
        \text{se e solo se}\notag\\ 
        \mathcal{A}^{I}/\mathcal{U}\vDash \phi\left[\pi(a^{0}),\pi(a^{1}), \dots, \pi(a^{n})\right]\label{eq:due}
    \end{gather}per qualche $ a^{0} \in A $.

    Per il teorema di Łoś, ricordando che $ \pi(a^{k}) = [\sequence{a^{k}_i}] $ in cui $ a^{k}_i \equiv a^{k}$ per ogni $ i \in I $, \eqref{eq:uno} è verificata se e solo se \begin{equation*}
        X'  =\set{i \in I: \mathcal{A}\vDash \exists y\,\phi\left[y,a^{1}_i, \dots, a^{n}_i\right]}  \in \mathcal{U}
    \end{equation*}Dal momento che $ a^{k}_i \equiv a^{k}$ per ogni $ i $, $ X' $ deve necessariamente essere $ I $ oppure $ \emptyset $. Facendo parte però di un ultrafiltro si ha che $ X'\neq \emptyset $, e pertanto $ X' = I $. 

    Segue dunque che \eqref{eq:uno} è equivalente a $ \mathcal{A}\vDash \exists y\,\phi\left[y,a^{1}, \dots, a^{n}\right] $, verificata se e solo se esiste $ a^0 \in A$ tale per cui \begin{equation*}
        \mathcal{A}\vDash \phi\left[a^{0},a^{1}, \dots, a^{n}\right]
    \end{equation*}Questo è equivalente a richiedere che\begin{equation*}
        X'' = \set{i \in I: \mathcal{A}_i\vDash \phi\left[a^{0},a^{1}_i,\dots,a^{n}_i\right]} = I
    \end{equation*}che è verificato, sempre per il teorema di Łoś, se e solo se\begin{equation*}
        \mathcal{A}^{I}/\mathcal{U}\vDash \phi\left[\pi(a^{0}),\pi(a^{1}), \dots, \pi(a^{n})\right].\qedd
    \end{equation*}
}
\osservazione{
    Il corollario di cui sopra definisce proprio il principio di transfer, ovvero l'immersione elementare di $ \R $ in $ \starR $.
}
\subsection{Iperreali come ultraprodotto}

\proposizione{dkdkjncjdjdkjncjskjncjdkjncjdkjncjskjndkjckjkjnsd}{
    Esiste un ultrafiltro non principale $ \mathcal{U} $ su $ \N $.
}
\dimostrazioneprop{dkdkjncjdjdkjncjskjncjdkjncjdkjncjskjndkjckjkjnsd}{
    Si vuole dimostrare che esiste un ultrafiltro $ \mathcal{U} \subseteq \parti{ \N} $ tale che \[
        A \subseteq \N\text{ finito} \,\implies\, \N \setminus A \in \mathcal{U}.
    \]

    Si consideri dunque il filtro di Fréchet su $ \N $: $ \mathcal{F}_{\text{Fr}}^{ \N} $. 
    
    Per il lemma \framref{doijcsoijoijoijoijoijoijoij} esiste $ \mathcal{U} $ ultrafiltro su $ \N $ tale che $ \mathcal{U}\supseteq \mathcal{F}_{\text{Fr}}^{ \N} $. Inoltre $ \mathcal{U} $ non può contenere alcun insieme finito. \qed
    
    % Infatti, se $ F \subseteq \N $ è finito e $ F \in \mathcal{U} $ 
    
    % $\implies$ $ \N\setminus F \in \mathcal{F}_{\text{Fr}}^{ \N}$ (per definizione di filtro di Fréchet) 
    
    % $\implies$ $ \N\setminus F \in \mathcal{U}$ poiché $ \mathcal{F}_{\text{Fr}}^{ \N} \subseteq \mathcal{U} $ 
    
    % $\implies$ $ F, \N\setminus F \in \mathcal{U} $ e quindi $ \emptyset \in \mathcal{U} $. Assurdo.\qed
}
\definizione{
    Si consideri il linguaggio dei campi ordinati: \[
        \mathcal{L}_{\text{ord}}=\{+,-, \cdot ,0,1,\le\}
    \]e la $ \mathcal{L}_{\text{ord}}$-struttura $ \mathcal{R} = (\R,+,-, \cdot ,0,1,\le ) $. 

    Si definisce il \emph{campo\footnote{Sarà chiaro tra qualche riga il perché lo si definisca un campo.} dei numeri reali non standard} $ \starR $ (anche \emph{campo degli iperreali}) come il dominio della $ \mathcal{L}_{\text{ord}} $-struttura \[
        \null^{*}\mathcal{R} \coloneqq \mathcal{R}^{ \N}/\mathcal{U}
    \]dove $ \mathcal{U} $ è un ultrafiltro non principale su $ \N $, ovvero \[
        \starR \coloneqq \R^{\N}/ \mathcal{U}.
    \]
}
\osservazione{
    Dalla costruzione segue che gli elementi di $ \starR $ siano classi di equivalenza (rispetto alla relazione indotta dall'ultrafiltro) di successioni di numeri reali $ \sequence{a_{i}}_{i \in \N} $. Si continuerà ad adottare la scrittura $ \sequence{a_{i}} $ per le successioni, sottintendendo che l'indice vari su $ \N $.
}
\paragrafo{Proprietà}{%
    Segue dal corollario \framref{dkjnckjjjjjjjjjjjjjdjjjjjdjnkjnkjn} che \begin{align*}
    \pi: \mathcal{R} &\longrightarrow  \null^{*}\mathcal{R}\\
    a &\longmapsto \left[\sequence{a}\right]
    \end{align*}è una immersione elementare, ed in particolare: \begin{itemize}
        \item $ \mathcal{R} $ e $ \null^{*} \mathcal{R} $ sono elementarmente equivalenti;
        \item questo garantisce che $ \starR $ sia un campo ordinato, poiché la proprietà ``essere un campo ordinato'' è esprimibile al prim'ordine\footnote{Si veda l'esempio 3.21 di \cite{andretta_elements_2022}};
        \item $ \R \subseteq \starR $ tramite l'immersione $ \pi $, che a ciascun numero reale $ a $ associa la classe di equivalenza della successione costante $ \sequence{a} $. Per indicare $ \pi(a) $ si utilizzerà la notazione $ \nonstandard{a} $.
    \end{itemize}
}{}{}
\osservazione{
    Segue immediatamente dalla costruzione di $ \starR $ che \begin{align*}
        0^{\null^{*}\mathcal{R}} &= \pi\left(0^{\mathcal{R}}\right) \eqqcolon \nonstandard{0};\\ 
        1^{\null^{*}\mathcal{R}} &= \pi\left(1^{\mathcal{R}}\right) \eqqcolon \nonstandard{1}.
    \end{align*}
}



\section{Gli iperreali}

In questa sezione si raccolgono alcune proprietà di $ \starR $, e si danno alcune definizioni.

\paragrafo{Osservazione}{%
    Facendo riferimento alla definizione di ultraprodotto fornita in \framref{dcjkklkmdjjjjultraprodottoinssimo}, somma, prodotto e relazione d'ordine su $ \starR $ sono definite come segue: per due elementi $ r, s $ di $ \starR $, $ r=\left[\sequence{r_{i}}\right]$, $s=\left[\sequence{s_{i}}\right]$\begin{align*}
        r+s &= \left[\sequence{r_{i}+s_{i}}\right];\\
        r \cdot s &= \left[\sequence{r_{i} \cdot s_{i}}\right];\\
        r\le s\quad &\text{se e solo se}\quad \{i \in \N: r_{i}\le s_{i}  \} \in \mathcal{U}.
    \end{align*}
}{}{}

\paragrafo{Osservazione}{%
    Per ogni $ r \in \starR $, $ r =\left[\sequence{r_{i}}\right] $, si definisce il valore assoluto come: \[
        |r| \coloneqq \left[\sequence{|r_{i}|}\right]
    \]Si noti che questa definizione agisce come ci si aspetta, ovvero:\begin{itemize}
        \item se $ r\ge \nonstandard{0} $ allora, per definizione, \begin{equation*}
            A\coloneqq \set{i \in \N: r_{i}\ge 0} \in \mathcal{U},
        \end{equation*}ma $ A = \set{i \in \N: r_{i} = |r_{i}|  } $, e pertanto $ \abs{r}=r $;
        \item Se $ r < \nonstandard{0} $ allora 
        \begin{equation*}
            B\coloneqq\{i \in \N: r_{i}< 0 \} \in \mathcal{U},
        \end{equation*}ma $ B= \{i \in \N: -r_{i} = |r_{i}|  \} $ e pertanto $ |r| = -r = \left[\sequence{-r_{i}}\right]$.
    \end{itemize}
}{}{}
\paragrafo{Notazione}{%
    Poiché è presente una inclusione naturale $ \R \hookrightarrow \starR $, non è necessario indicare con simboli diversi $ r \in R $ e $ \nonstandard{r} = \pi(r) \in \starR $. 

    Pertanto, per tutti gli elementi di $ r \in \R $, quando questo non creerà ambiguità si indicherà con $ r $ anche $ \nonstandard{r} $. 

    La stessa cosa si applica a relazioni e operazioni: quando ciò non crea ambiguità, si eviterà di indicare ``$ \nonstandard{\null} $''.
}{}{}
\proposizione{djkjncdjcjndkjncjdjckjncjjdkjncjdkjncsdkjncskjndjncdksj}{
    $ \starR $ è un campo non archimedeo.
}
\dimostrazioneprop{djkjncdjcjndkjncjdjckjncjjdkjncjdkjncsdkjncskjndjncdksj}{
    
    % Dimostrando l'esistenza di $ M \in \starR  $ tale per cui \[
    %     M \ge n1\qquad \forall\, n \in \N
    % \]si Dunque si ha che \[
    %     n1 = \left[\sequence{n}\right]
    % \]

    Sia $ \sequence{a_{i}}_{i \in \N} $ la successione tale che per ogni $ i$, $a_{i} = i$ e sia $ M \coloneqq \left[\sequence{a_{i}}\right] $. \[
        M=\left[\sequence{0,1,2,3,\dots}\right]
    \]

    Fissato ora $ n \in \N $, $ n1 \le M $ se e solo se l'insieme $ \displaystyle A\coloneqq\{i \in \N: n \le   i\} $ appartiene ad $ \mathcal{U}$. 

    Sia ora $ B \coloneqq \N\setminus A $ (e quindi $ A = \N\setminus B $): \[
        B = \{0,1,2,\dots,n-1\}
    \]ed è un insieme \emph{finito}; segue che $ A \in \mathcal{U}$ poiché si è scelto $ \mathcal{U} $ ultrafiltro non principale.
    
    Pertanto $ M $ è un elemento di $ \starR $, ma per ciascun $ n \in \N $ si ha che $ M\ge n1 $: $ \starR $ non è un campo archimedeo.\qed
}

\definizione{
    Un numero iperreale $ x \in \starR $ si dice \begin{itemize}
        \item \emph{finito} se esiste un numero reale $ r \in \R $ tale per cui $ |x|<r $;
        \item \emph{infinito} se per ogni numero reale $ r \in \R$, $ r<|x| $;
        \item \emph{infinitesimo} se per ogni numero reale positivo $ r \in \R $, $ r>0 $, vale \begin{equation*}
            0 < |x| < r.
        \end{equation*}
    \end{itemize}
}

\definizione{
    Sia $ X \subseteq \R $. Il sottoinsieme $ \nonstandard{X} \subseteq \starR $ si definisce come \begin{equation*}
        \nonstandard{X} \coloneqq \set{\left[\sequence{a_{n}}\right]: a_{n} \in X, \forall\, n \in \N}.
    \end{equation*}
}
\definizione{
    Sia $ X \subseteq \R $ e sia $ f: X \longrightarrow \R $. Si definisce l'estensione non standard di $ f $ come \begin{align*}
    \nonstandard{f}:\nonstandard{X} &\longrightarrow \starR \\
    \left[\sequence{a_{n}}\right] &\longmapsto \left[\sequence{f(u_{n} )}\right].
    \end{align*}
}

\paragrafo{Osservazione}{%
    Date le definizioni di cui sopra, è possibile ampliare il principio di transfer. 

    In particolare se $ \phi(x_1,\dots,x_{n} ) $ è una formula che contiene unicamente elementi di $ \mathcal{L}_{ord} $ oppure espressioni come $ f(x) $ e $ x \in D $, dove $ f:\R \longrightarrow \R $ è una funzione e $ D \subseteq \R $ è un sottoinsieme specifico di $ \R $, allora $ \phi[a_1,\dots,a_{n} ] $ è verificata su $ \R $ se e solo se $ \tilde{\phi}[\nonstandard{a_1},\dots,\nonstandard{a_{n} }] $ è verificata su $ \nonstandard{\R} $, dove $ \tilde{\phi} $ è la formula ottenuta sostituendo a ciascun simbolo di funzione $ f $ la sua estensione non standard $ \nonstandard{f} $, e sostituendo a ciascun simbolo di sottoinsieme $ D $ la sua estensione non standard. 

    La dimostrazione di ciò sarà l'argomento del capitolo successivo.
}{tantebellecoseinoijknkjnkjnsd}{}

\subsection{Elementi interni}

\definizione{
    Un elemento $\xi \in \starR$ si dice interno se appartiene a qualche iper-estensione $ \nonstandard{\Upsilon } $ per $ \Upsilon  \subseteq \starR $. Un elemento dell’universo non standard che non è interno si dice esterno.
}

Si riporta una versione più debole del teorema 2.8.12 di \cite{loeb_nonstandard_2015}, la cui dimostrazione prescinde dall'obiettivo di questa tesi. 

\teorema{teoriemidkjnkjnc}{
    Sia $ A $ un sottoinsieme interno di $ \nonstandard{R} $. \begin{enumerate}
        \item Se $ \N \subseteq A $, allora contiene almeno un elemento di $ \nonstandard{\N} $ infinito. (Principio dell'\straniero{overspill}).
        \item Se $ A $ contiene tutti gli elementi infiniti di $ \nonstandard{\N} $, allora contiene almeno un elemento finito di $ \N $. (Principio dell'\straniero{underspill}).
    \end{enumerate}
}

\subsection{Numeri naturali e loro estensione}

\teorema{dkjncjskjnskjncjdjckjnsjdkjnckjnskjnd}{
    Ogni numero in $ \nonstandard{\N} $ finito è un numero naturale
}
\dimostrazione{dkjncjskjnskjncjdjckjnsjdkjnckjnskjnd}{
    Sia $ N \in \N $ fissato. I numeri naturali soddisfano la seguente formula: \begin{equation*}
        ( \forall\,n \in\N)\left[n<N \,\implies\,n = 0 \,\lor\,n=1\,\lor\, \dots \,\lor\, n = N\right]
    \end{equation*}Applicando il principio di transfer, si ottiene che \begin{equation*}
        ( \forall\,n \in\nonstandard{\N})\left[n<N \,\implies\,n = 0 \,\lor\,n=1\,\lor\, \dots \,\lor\, n = N\right]
    \end{equation*}Per l'arbitrarietà di $ N $, questo implica che se un numero ipernaturale $ \xi $ è finito, allora è necessariamente un numero naturale.\qed
}


È immediato dunque questo corollario:
\paragrafo{Corollario}{%
    Ogni elemento di $ \nonstandard{\N} $ è un numero naturale oppure un numero infinito.
}{}{}